% Part: computability
% Chapter: recursive-functions
% Section: primes

\documentclass[../../../include/open-logic-section]{subfiles}

\begin{document}

% SEGMENT OLP-0219-S01

\olfileid{cmp}{rec}{pri}
\olsection{பகா எண்கள்}

இயல்பான சார்புகளும் உறவுகளும் முதன்மை மீள்வரையறைக்குரியவை என்று
காட்டுவதற்குப் பெருமளவு கருவிகளை வரம்பிட்ட அளவையிடலும் வரம்பிட்ட
சிறுமமாக்கலும் வழங்குகின்றன. எடுத்துக்காட்டாக, “$x$ ஆனது $y$ ஐ
வகுக்கிறது” என்ற உறவைக் கருதுக; இதை $x \mid y$ என எழுதுகிறோம்.
மீதியின்றி $y$ ஐ~$x$ ஆல் வகுக்க முடிந்தால், அதாவது $y$ ஆனது
~$x$ இன் ஒரு முழுஎண் மடங்காக இருந்தால், $x \mid y$ பொருந்தும்.
(அது பொருந்தாவிட்டால், அதாவது $x$ ஐ $y$ ஆல் வகுக்கும்போது கிடைக்கும்
மீதி $> 0$ எனில், $x \nmid y$ என எழுதுகிறோம்.) வேறுவிதமாகக்
கூறினால், ஏதேனும்~$z$ க்கு $x \cdot z = y$ எனில், எனில் மட்டுமே,
$x \mid y$. அத்தகைய $z$ இருந்தால் அது $\leq y$ ஆகவே இருக்க வேண்டும்
என்பது தெளிவு. எனவே ஏதேனும் $z \le y$ க்கு $x \cdot z = y$ எனில்,
எனில் மட்டுமே, $x \mid y$. பின்வருமாறு $=$ மற்றும் பெருக்கலிலிருந்து
வரம்பிட்ட உளதளவையிடலால் $x \mid y$ என்ற உறவை வரையறுக்கலாம்:
\[
x \mid y \defiff \bexists{z \leq y}{(x \cdot z) = y}.
\]
இவ்வாறு $x \mid y$ ஒரு முதன்மை மீள்வரையறை உறவு என்று காட்டியுள்ளோம்.

இயல் எண்~$x$, அது $0$ உம் $1$ உம் அல்லாமல், $1$ ஆலும் தன்னாலும்
மட்டுமே வகுபட்டால் \emph{பகா எண்} ஆகும். வேறுவிதமாகக் கூறினால்,
$y \mid x$ பொருந்தும்போதெல்லாம் $y = 1$ அல்லது~$y=x$ என அமையும்
எண்களே பகா எண்கள். $x$ பகா எண்ணா என்று சோதிக்க, அனைத்து
$y \le x$ க்கும் $y \mid x$ பொருந்துகிறதா என்று மட்டும் பார்க்க
வேண்டும்; ஏனெனில் $y > x$ எனில் தானாகவே~$y \nmid x$.
எனவே $x$ பகா எண்ணாக இருந்தால், எனில் மட்டுமே, பொருந்தும் உறவு
$\fn{Prime}(x)$ ஐப் பின்வருமாறு வரையறுக்கலாம்:
\[
\fn{Prime}(x) \defiff x \geq 2 \land \bforall{y \leq x}{(y \mid x \lif y
  = 1 \lor y = x)}
\]
ஆகவே அது ஒரு முதன்மை மீள்வரையறை உறவு.

பகா எண்கள் $2$, $3$, $5$, $7$, $11$, முதலியவை. இத்தொடரில்
$x$-ஆவதாக வரும் பகா எண்ணைத் திருப்பித்தரும் சார்பு $p(x)$ ஐக்
கருதுக; அதாவது $p(0) = 2$, $p(1) = 3$, $p(2) = 5$, முதலியவை.
(வசதிக்காக $p(x)$ ஐப் பலவேளைகளில் $p_x$ என எழுதுவோம்
($p_0=2$, $p_1=3$, முதலியவை.)

~$x$ ஐவிடப் பெரிய முதல் பகா எண்ணைத் திருப்பித்தரும் சார்பு
$\fn{nextPrime(x)}$ நம்மிடம் இருந்தால், முதன்மை மீள்வரையறையைப்
பயன்படுத்தி~$p$ ஐ எளிதாக வரையறுக்கலாம்:
\begin{align*}
  p(0) & = 2\\
  p(x+1) & = \fn{nextPrime}(p(x))
\end{align*}
$\fn{nextPrime}(x)$ என்பது $y > x$ ஆகவும்~$y$ பகா எண்ணாகவும்
இருக்கும் மிகச்சிறிய $y$ என்பதால், வரம்பற்ற தேடலால் அதை எளிதாகக்
கணிக்கலாம். ஆனால் யூக்ளிடின் ஒரு முடிவினால், வரம்பிட்ட சிறுமமாக்கலைப்
பயன்படுத்தியும் அதை வரையறுக்கலாம்: $x$ க்கும் $\fact{x}+1$ க்கும்
இடையில் எப்போதும் ஒரு பகா எண் உள்ளது.
\[
  \fn{nextPrime(x)} =
  \bmin{y \leq \fact{x}+1}{(y > x \land \fn{Prime}(y))}.
\]
இதிலிருந்து $\fn{nextPrime}(x)$ உம், எனவே $p(x)$ உம்
(கணிக்கத்தக்கவை மட்டுமன்றி) முதன்மை மீள்வரையறைச் சார்புகள் என்று
தெரிகிறது.

(ஆர்வம் இருந்தால், யூக்ளிடின் தேற்றத்திற்கான விரைவான நிரூபிப்பு இதோ.
$p_n$ ஆனது $\le x$ ஆக உள்ள மிகப்பெரிய பகா எண் என்றும்,
$p = p_0\cdot p_1 \cdot \dots \cdot p_n$ ஆனது~$\le x$ ஆக உள்ள
அனைத்துப் பகா எண்களின் பெருக்குத்தொகை என்றும் கொள்க. $p+1$ ஒரு
பகா எண்ணாக இருக்கும், அல்லது $x$ க்கும்~$p+1$ க்கும் இடையில்
ஒரு பகா எண் இருக்கும். ஏன்? $p+1$ ஒரு பகா எண் அல்ல என்று கொள்க.
அப்போது $q < p+1$ ஆக உள்ள ஏதேனும் பகா எண் $q \mid p+1$.
$\le x$ ஆக உள்ள பகா எண்களுள் எதுவும் $p+1$ ஐ வகுக்காது.
(~$p$ இன் வரையறைப்படி, $p_i \le x$ ஆக உள்ள ஒவ்வொரு பகா எண்ணும்
~$p$ ஐ வகுக்கிறது; அதாவது மீதி~$0$. எனவே $p_i \le x$ ஆக உள்ள
ஒவ்வொரு பகா எண்ணாலும் $p+1$ ஐ வகுக்கும்போது மீதி~$1$; ஆகவே
$p_i \nmid p+1$.) இதனால் $q$ என்பது $> x$ மற்றும் $< p+1$
ஆக உள்ள ஒரு பகா எண். மேலும் $p \le \fact{x}$; எனவே $> x$ மற்றும்
$\le \fact{x}+1$ ஆக உள்ள ஒரு பகா எண் இருக்கிறது.)

\begin{prob}
வரம்பிட்ட சிறுமமாக்கலைப் பயன்படுத்தி முழுஎண் வகுத்தல் $d(x, y)$ ஐ
வரையறுக்க.
\end{prob}

\end{document}
